\documentclass[a4paper,reqno,10pt]{amsart}

\usepackage{amsfonts,amssymb,amsthm,amsmath,amscd,mathtools,mathrsfs}
\usepackage{graphicx,color,xcolor,cite,leftidx,comment,enumitem}
\usepackage{fancybox,multirow,makecell,caption,subcaption}
\usepackage{tasks,thmtools,array,booktabs}
\usepackage[all,2cell,cmtip]{xy}
\xyoption{curve}
\usepackage{tikz,tikz-cd,float}
\usetikzlibrary{arrows,decorations.pathmorphing,decorations.markings,
backgrounds,positioning,fit,petri,patterns,matrix,shapes.geometric}
\usepackage[hypertexnames=false]{hyperref}
\usepackage[nameinlink]{cleveref}

\oddsidemargin=10pt
\evensidemargin=6pt
\textwidth=15.5true cm
\textheight=21.6true cm
\renewcommand{\baselinestretch}{1.20}
\setlength{\parskip}{5pt}
\let\shorttwoheadrightarrow\twoheadrightarrow
\renewcommand{\hookrightarrow}{\lhook\joinrel\longrightarrow}
\renewcommand{\twoheadrightarrow}{\relbar\joinrel\relbar\joinrel
\shorttwoheadrightarrow}

\hypersetup{
  colorlinks=true,
  linkcolor=blue!55!black,
  citecolor=green!45!black,
  urlcolor=blue!60!black,
  pdftitle={A pre-triangulated category which is not triangulated},
  pdfauthor={Xiao-Wu Chen, Jian Liu, Xue-Song Lu, Chencheng Zhang}
}

\declaretheoremstyle[headfont=\bfseries,bodyfont=\itshape]{plainstyle}
\declaretheoremstyle[headfont=\bfseries,bodyfont=\normalfont]{defstyle}
\declaretheoremstyle[headfont=\bfseries,bodyfont=\normalfont]{remarkstyle}
\declaretheorem[style=plainstyle,numberwithin=section,name=Theorem]{theorem}
\declaretheorem[style=plainstyle,sibling=theorem,name=Lemma]{lemma}
\declaretheorem[style=plainstyle,sibling=theorem,name=Proposition]{proposition}
\declaretheorem[style=plainstyle,sibling=theorem,name=Corollary]{corollary}
\declaretheorem[style=plainstyle,numbered=no,name=Theorem]{restatement}
\declaretheorem[style=defstyle,sibling=theorem,name=Definition]{definition}
\declaretheorem[style=defstyle,sibling=theorem,name=Convention]{convention}
\declaretheorem[style=plainstyle,sibling=theorem,name=Conjecture]{conjecture}
\declaretheorem[style=remarkstyle,sibling=theorem,name=Remark]{remark}
\numberwithin{equation}{section}


\newtheorem*{ack}{Acknowledgements}

\DeclareMathOperator{\modu}{mod}
\DeclareMathOperator{\End}{End}
\DeclareMathOperator{\Hom}{Hom}
\newcommand{\kk}{\mathbb F_2}
\newcommand{\cF}{\mathcal F}
\newcommand{\cO}{\mathcal O}
\newcommand{\Lam}{\Lambda}
\newcommand{\eps}{\epsilon}
\newcommand{\ol}[1]{\underline{#1}}
\newcommand{\longhookrightarrow}{\lhook\joinrel\longrightarrow}
\setcounter{MaxMatrixCols}{12}

\title[A pre-triangulated category which is not triangulated]
{A pre-triangulated category which is not triangulated}

\author{Xiao-Wu Chen}
\address{School of Mathematical Sciences, University of Science and
Technology of China, Hefei 230026, Anhui Province, P. R. China}
\email{xwchen@mail.ustc.edu.cn}

\author{Jian Liu}
\address{School of Mathematics and Statistics, and Hubei Key Laboratory of Mathematical Sciences, Central China Normal University, Wuhan 430079, P. R. China}
\email{jianliu@ccnu.edu.cn}

\author{Xue-Song Lu}
\address{School of Mathematical Sciences, Shanghai Jiao Tong University,
Shanghai 200240, P. R. China}
\email{leocedar@sjtu.edu.cn}

\author{Chencheng Zhang}
\address{School of Mathematical Sciences, Shanghai Jiao Tong University,
Shanghai 200240, P. R. China}
\email{zhangchencheng@sjtu.edu.cn}

\subjclass[2020]{Primary 18G80; Secondary 16G20, 16G70}
\keywords{pre-triangulated category, triangulated category, octahedral axiom, Heller comparison, preprojective algebra}

\begin{document}

\begin{abstract}
In this article, we construct an explicit pre-triangulated category which is not a triangulated category.  Its underlying additive category is the category of finitely generated projective modules of the type-$A_5$ preprojective algebra over $\mathbb F_2$, and the suspension is induced by the graph-reflection automorphism. 
\end{abstract}

\maketitle

\section{Introduction}

In the 1960s, Grothendieck introduced derived categories; see \cite{Hartshorne1966}. Verdier \cite{Verdier1977} subsequently developed the theory of \emph{triangulated categories} and formulated the four axioms $\mathbf{TR1}$--$\mathbf{TR4}$ that characterize their structure.
Since then, triangulated categories have become a fundamental framework in homological algebra,
algebraic geometry, and representation theory. Classical examples of triangulated categories
include derived categories and stable categories of Frobenius categories; see
\cite{Happel1988, Neeman2001}.
Among these axioms, the fourth one, the octahedral axiom $\mathbf{TR4}$, controls the interaction
between composition of morphisms and mapping cones.
If only $\mathbf{TR1}$--$\mathbf{TR3}$ are required, the resulting
structure is called \emph{pre-triangulated}. A pre-triangulated category is also called Puppe-triangulated; see \cite{Neeman2001, Puppe1967}.


Let $\mathcal B$ be an abelian Frobenius category. That is, $\mathcal B$ is an abelian category with enough projective objects and enough injective objects, and the projective objects coincide with the injective objects. A classical result of Happel \cite{Happel1988} shows that the stable category of $\mathcal B$ modulo projective objects, denoted by $\underline{\mathcal B}$, admits a triangulated structure. Its suspension functor is given by the cosyzygy functor
$\Omega^{-1}\colon \underline{\mathcal B}\xrightarrow \cong \underline{\mathcal B}$.
The category $\mathcal B$ is called \emph{$\Sigma$-stable} in the sense of Beligiannis \cite{Beligiannis2000} if there exists a stable structure $(\Sigma,\delta)$. Here, $(\Sigma,\delta)$ is a stable structure provided that $\Sigma\colon \mathcal B\to\mathcal B$ is an auto-equivalence and there exists a natural isomorphism
$\delta\colon \Sigma\xrightarrow \cong \Omega^{-3}$
in $\underline{\mathcal B}$ such that
$\delta_{\Omega^{-1} N}=-\Omega^{-1}(\delta_N)\circ \sigma_N$ for each $N\in \mathcal B$, where $\sigma: \Sigma\Omega^{-1}\xrightarrow \cong \Omega^{-1}\Sigma$ is the natural isomorphism induced by $\Sigma$.

Given an abelian Frobenius category $\mathcal B$ with an auto-equivalence 
$\Sigma\colon \mathcal B\xrightarrow\cong \mathcal B$, a $\Sigma$-stable structure $(\Sigma,\delta)$ on $\mathcal B$ 
gives rise to a pre-triangulated category $(\mathcal P,\Sigma,\triangle_\delta)$ constructed by Heller \cite{Heller1968}. 
Here, $\triangle_\delta$ consists of all distinguished triangles in this category, and 
$\mathcal P$ denotes the full subcategory of projective objects of $\mathcal B$; see \Cref{thm:heller}. 
Moreover, Heller \cite[Theorem~16.4]{Heller1968} showed that this construction establishes a one-to-one correspondence between $\Sigma$-stable structures on $\mathcal B$ and pre-triangulated category structures on $(\mathcal P,\Sigma)$. 
Amiot found a sufficient condition ensuring that the pre-triangulated structure constructed by Heller's method is triangulated, and applied this result to deformed preprojective algebras; see \cite[Section 9]{Ami07}.



It is natural to ask whether every pre-triangulated category must be triangulated. 
This question remained open for a long time. 
In \cite[Conjecture~9.9]{Beligiannis2000}, Beligiannis formulated the following conjecture, which he attributed to Keller and Neeman.

\begin{conjecture}
There exists an abelian Frobenius category $\mathcal B$ with a $\Sigma$-stable structure $(\Sigma,\delta)$ such that the induced pre-triangulated category $(\mathcal P,\Sigma,\triangle_\delta)$ is not triangulated. Here, $\mathcal P$ is the full subcategory of $\mathcal B$ consisting of all projective objects.
\end{conjecture}

The main result of this article is \Cref{thm:main}, which confirms the above conjecture. It shows that $\mathbf{TR4}$ is not
a formal consequence of $\mathbf{TR1}$--$\mathbf{TR3}$ in general. 


 In this article, we study the preprojective algebra $\Pi(A_5)$ of the Dynkin quiver of type $A_5$ over $\mathbb F_2$. 
This algebra is self-injective. In particular, the category $\modu \Pi(A_5)$ of finitely generated right modules is an abelian Frobenius category. 
Thus, its stable category $\ol{\modu}\Pi(A_5)$ is a triangulated category. 
Moreover, the category $\mathcal P(\Pi(A_5))$ of finitely generated projective $\Pi(A_5)$-modules admits a triangulated category structure; see \Cref{conv:transported-triangulation}. 
The corresponding $\Sigma$-stable structure on $\modu \Pi(A_5)$ is given by the pair $(\Sigma, \theta_0)$, where $\theta_0:\Sigma\xrightarrow{\sim}\Omega^{-3}$ is an isomorphism on $\ol{\modu}\Pi(A_5)$; here, $\Sigma$ is induced by the graph-reflection automorphism; see \Cref{conv:transported-triangulation}.


By choosing a module $M$ whose second syzygy is isomorphic
to itself, we construct a twisted comparison
$$\theta_\eps:\Sigma\xrightarrow{\sim}\Omega^{-3}$$ 
such that $(\Sigma, \theta_\epsilon)$ is a stable structure; see \Cref{prop:twisted-comparison}.
Thus Heller's theorem \cite[Theorem 16.4]{Heller1968} yields a pre-triangulated category
$(\mathcal P(\Pi(A_5)), \Sigma, \triangle_{\theta_\eps})$ satisfying $\mathbf{TR1}$--$\mathbf{TR3}$. The following main result in this article shows that
$\mathbf{TR4}$ fails for $\triangle_\eps$.

\begin{theorem}\label{thm:main}
The abelian Frobenius category $\modu \Pi(A_5)$ admits a $\Sigma$-stable structure $(\Sigma, \theta_\eps)$ such that the induced pre-triangulated category $(\mathcal P(\Pi(A_5)), \Sigma, \triangle_{\theta_\eps})$ is not triangulated. 
\end{theorem}

$\triangle_{\theta_{\eps}}$ is simply written as $\triangle_{\eps}$ in the main text of the article. The proof of the above result is given at the end of
\Cref{sec:example}. It exhibits two distinguished rows and two distinguished
columns which admit no four-by-four completion. The key point is a relative
rigidity statement: every possible third vertical arrow can be normalized by
an automorphism of the second row, after which the third row is forced to have
the original Heller comparison rather than the twisted one.


\section{Preliminaries}\label{sec:preliminaries}

\subsection{(Pre)triangulated categories}\label{sec:heller}

Throughout, $k=\mathbb F_2$. All modules are finite-dimensional right modules.
For a finite-dimensional algebra $A$, let $\mathcal P(A)$ denote the full
subcategory of finitely generated projective right $A$-modules.

\begin{definition}\label{def:stable-category}
If $X$ and $Y$ are $A$-modules, set
\[
 \mathcal P(X,Y)=
 \{g\circ f\mid
 X\xlongrightarrow{f}P\xlongrightarrow{g}Y
 \ \text{ for some projective module } P\}.
\]
That is, $\mathcal P(X,Y)$ is the subspace of morphisms which factor through a
projective module, and the \emph{stable Hom-space} is
\[
 \ol{\Hom}_A(X,Y)
 =\Hom_A(X,Y)/\mathcal P(X,Y).
\]
The \emph{stable category} has the same objects and these stable Hom-spaces. For a Frobenius category $\mathfrak B$,
we write $\Omega$ for the syzygy autoequivalence of its stable category and
$S=\Omega^{-1}$ for the stable suspension.
\end{definition}

\begin{definition}\label{def:verdier-axioms}
Let $\mathcal T$ be an additive category. An \emph{autoequivalence} of
$\mathcal T$ is an additive functor
$\Sigma:\mathcal T\longrightarrow\mathcal T$ with an additive
quasi-inverse; we call it the \emph{suspension}. A
\emph{candidate $3$-$\Sigma$-sequence} is a sequence
\[
\xymatrix@C=38pt{
X\ar[r]^-{u}&Y\ar[r]^-{v}&Z\ar[r]^-{w}&\Sigma X.
}
\]
No vanishing condition is imposed in this definition. Its
\emph{forward rotation} is
\[
\xymatrix@C=38pt{
Y\ar[r]^-{v}&Z\ar[r]^-{w}&\Sigma X\ar[r]^-{-\Sigma u}&\Sigma Y.
}
\]
Let $\triangle$ be a class of candidate $3$-$\Sigma$-sequences; its
members are called \emph{distinguished}. We use Verdier's original axioms
\cite[Definition~1-1]{Verdier1977}; see also \cite[\S1.1]{Neeman2001}.

\begin{description}[leftmargin=3.6em,style=nextline]
\item[$\mathbf{TR1}$]
The class $\triangle$ is closed under isomorphisms. For every morphism $u:X\longrightarrow Y$, there are
$Z,v,w$ for which the resulting candidate sequence is distinguished. For
every object $X$, the sequence
\[
\xymatrix@C=38pt{
X\ar[r]^-{\mathrm{Id}_X}&X\ar[r]&0\ar[r]&\Sigma X
}
\]
is distinguished.

\item[$\mathbf{TR2}$]
A candidate sequence is distinguished if and only if its forward
rotation is distinguished.

\item[$\mathbf{TR3}$]
Given distinguished rows and a commutative square on their first maps, a
third vertical map exists which makes the entire diagram commute:
\[
\xymatrix@C=35pt@R=15pt{
X\ar[r]^-{u}\ar[d]_-{a}&
Y\ar[r]^-{v}\ar[d]_-{b}&
Z\ar[r]^-{w}\ar@{-->}[d]^-{c}&
\Sigma X\ar[d]^-{\Sigma a}\\
X'\ar[r]_-{u'}&Y'\ar[r]_-{v'}&Z'\ar[r]_-{w'}&\Sigma X'.
}
\]
Thus $b\circ u=u'\circ a$, $c\circ v=v'\circ b$ and $w'\circ c=\Sigma a\circ w$.
The dotted map need not be unique. The resulting triple $(a,b,c)$ is a
\emph{morphism of candidate $3$-$\Sigma$-sequences}. 

\item[$\mathbf{TR4}$]
For any distinguished triangles
$X\xlongrightarrow{u}Y\xlongrightarrow{i}X'
\xlongrightarrow{i'}\Sigma X$,
$Y\xlongrightarrow{v}Z\xlongrightarrow{j}Z'
\xlongrightarrow{j'}\Sigma Y$, and
$X\xlongrightarrow{v\circ u}Z\xlongrightarrow{k}Y'
\xlongrightarrow{k'}\Sigma X$, there are morphisms
$u':X'\longrightarrow Y'$ and $v':Y'\longrightarrow Z'$ such that the
following diagram commutes:
\[
\xymatrix@C=30pt@R=15pt{
X\ar[r]^-u\ar@{=}[d]&
Y\ar[r]^-i\ar[d]^-v&
X'\ar[r]^-{i'}\ar[d]^-{u'}&
\Sigma X\ar@{=}[d]\\
X\ar[r]^-{v\circ u}&
Z\ar[r]^-k\ar[d]_-j&
Y'\ar[r]^-{k'}\ar[d]^-{v'}&
\Sigma X\ar[d]^-{\Sigma u}\\
&Z'\ar@{=}[r]\ar[d]_-{j'}&
Z'\ar[r]^-{j'}\ar[d]^-{(\Sigma i)\circ j'}&
\Sigma Y\\
&\Sigma Y\ar[r]_-{\Sigma i}&\Sigma X'.
}
\]
Moreover, the sequence
$X'\xlongrightarrow{u'}Y'\xlongrightarrow{v'}Z'
\xlongrightarrow{(\Sigma i)\circ j'}\Sigma X'$ is distinguished.
This is Verdier's original \emph{octahedral axiom}, unfolded into the plane.
\end{description}

Recall that $(\mathcal T,\Sigma,\triangle)$ is \emph{pre-triangulated} if it
satisfies $\mathbf{TR1}$--$\mathbf{TR3}$, and \emph{triangulated} if it also satisfies
$\mathbf{TR4}$; see \cite[\S9.1]{Beligiannis2000}. 
For a Frobenius category $\mathfrak B$,  a classical result of Happel \cite[Theorem 2.6]{Happel1988} shows that the stable category of $\mathfrak B$ with suspension functor $\Omega^{-1}$ can form a triangulated category. 
\end{definition}

\begin{remark}\label{rem:consecutive-zero}
Under
$\mathbf{TR1}$--$\mathbf{TR3}$, the composition of any two consecutive morphisms in a distinguished triangle is zero; see  \cite[Remark~1.1.3]{Neeman2001}. Namely, if $(\mathcal T,\Sigma,\triangle)$ is a pre-triangulated category, then for each distinguished triangle $
\xymatrix@C=38pt{
X\ar[r]^-{u}&Y\ar[r]^-{v}&Z\ar[r]^-{w}&\Sigma X
}
$
in $\triangle$, one has $w\circ v=0, v\circ u=0$ and $(\Sigma u)\circ w=0$. 
\end{remark}

\begin{definition}\label{def:fixed-boundary}
A pre-triangulated category has the \emph{four-by-four property} if every choice of distinguished first two rows and first two
columns in the diagram
\[
\xymatrix@C=27pt@R=15pt{
A_0\ar[r]^-{a_0}\ar[d]_-{x_0}&
A_1\ar[r]^-{a_1}\ar[d]_-{x_1}&
A_2\ar[r]^-{a_2}\ar@{-->}[d]^-{x_2}&
\Sigma A_0\ar[d]^-{\Sigma x_0}\\
B_0\ar[r]^-{b_0}\ar[d]_-{y_0}&
B_1\ar[r]^-{b_1}\ar[d]_-{y_1}&
B_2\ar[r]^-{b_2}\ar@{-->}[d]^-{y_2}&
\Sigma B_0\ar[d]^-{\Sigma y_0}\\
C_0\ar@{-->}[r]^-{c_0}\ar[d]_-{z_0}&
C_1\ar@{-->}[r]^-{c_1}\ar[d]_-{z_1}&
C_2\ar@{-->}[r]^-{c_2}\ar@{-->}[d]^-{z_2}&
\Sigma C_0\ar[d]^-{-\Sigma z_0}\\
\Sigma A_0\ar[r]_-{\Sigma a_0}&
\Sigma A_1\ar[r]_-{\Sigma a_1}&
\Sigma A_2\ar[r]_-{-\Sigma a_2}&
\Sigma^2 A_0
}
\]
whose north-west square commutes can be completed so that all rows and
columns are distinguished and all small squares commute, except for the
south-east square, which anticommutes.
\end{definition}

\begin{lemma}\label{lem:tr4-implies-fixed-boundary}
Every triangulated category has the four-by-four property.
\end{lemma}

\begin{proof}
This is the three-by-three lemma, which follows from the octahedral axiom;
see \cite[Lemma~2.6]{May2001}. We use only this implication.
\end{proof}

\begin{convention}\label{conv:heller-coherence}
Let $A$ be a finite-dimensional self-injective algebra, set
$\cF=\mathcal P(A)$, and let $\Sigma$ be an exact autoequivalence of
$\modu A$ which restricts to $\cF$. Write $S=\Omega^{-1}$ for the
suspension of $\ol{\modu}A$. Fixed projective-injective resolutions give a
natural isomorphism
\[
 \sigma_N:\Sigma(SN)\xlongrightarrow{\sim}S(\Sigma N).
\]
Thus $(\Sigma,\sigma)$ is a triangle functor. A natural isomorphism
$\theta:\Sigma\longrightarrow S^3$ is an isomorphism of triangle functors
precisely when
\begin{equation}\label{eq:triangle-functor-compatibility}
 \theta_{SN}=-S(\theta_N)\circ\sigma_N;
\end{equation}
see \cite[Theorem~16.4]{Heller1968}, the sign convention here is from the forward rotation.
\end{convention}

\begin{definition}
\label{def:complete-comparison}
A candidate $3$-$\Sigma$-sequence in $\cF$ is \emph{exact} if its
$\Sigma$-periodic continuation is an exact sequence of $A$-modules. In
particular, all consecutive composites vanish.

Denote an exact candidate
$3$-$\Sigma$-sequence by
\[
 \mathfrak s=
 \bigl(X\xlongrightarrow{u}Y\xlongrightarrow{v}Z
 \xlongrightarrow{w}\Sigma X\bigr),
 \qquad N_{\mathfrak s}=\operatorname{Ker}u.
\]
The sequence $\mathfrak s$ begins an injective resolution of
$N_{\mathfrak s}$. Comparing it with a fixed injective resolution yields a
stable isomorphism
\[
  \delta_{\mathfrak s}:\Sigma N_{\mathfrak s}
  \longrightarrow\Omega^{-3}(N_{\mathfrak s}).
\]
We call this the \emph{complete Heller comparison}. 
\end{definition}

The following classification is due to Heller; its formulation by exact
$3$-$\Sigma$-sequences is due to Geiss--Keller--Oppermann.

\begin{theorem}[{\cite[Theorem~16.4]{Heller1968};
\cite[Lemma~2.3 and Proposition~2.4]{GeissKellerOppermann2013}}]
\label{thm:heller}
In the situation of \Cref{conv:heller-coherence}, for every isomorphism of
triangle functors
\[
 \theta:(\Sigma,\sigma)\xlongrightarrow{\sim}
(\Omega^{-3},-\mathrm{Id}_{\Omega^{-4}}),
\]
the exact candidate $3$-$\Sigma$-sequences $\mathfrak s$ satisfying $\delta_{\mathfrak s}=\theta_{N_{\mathfrak s}}$
form a class $\triangle_\theta$ and make
$(\cF,\Sigma,\triangle_\theta)$ pre-triangulated. Every such structure on
$(\cF,\Sigma)$ arises uniquely in this way.
\end{theorem}

\begin{proof}
Because $A$ is a projective generator, evaluation at $A$ identifies the
exactness occurring in Heller's construction for $\mathcal P(A)$ with
exactness of $A$-modules. Then we could apply
\cite[Theorem~16.4]{Heller1968} to the present setting. The
formulation by exact $3$-$\Sigma$-sequences and complete comparisons is the
case $n=3$ of
\cite[Lemma~2.3 and Proposition~2.4]{GeissKellerOppermann2013}. The
compatibility with $\Omega^{-1}$ is exactly the condition that $\theta$ be
an isomorphism of triangle functors. In characteristic two the displayed
minus sign equals the identity, but the triangle-functor condition remains
part of the hypothesis.
\end{proof}


Applying $\Omega^3$ to $\theta_N$ and using the standard isomorphism
$\Omega^3\circ\Omega^{-3}\cong
\mathrm{Id}_{\ol{\modu}A}$, we get
\[
  \lambda_N=\Omega^3(\theta_N):\Omega^3(\Sigma N)\longrightarrow N.
\]
 For an exact candidate $3$-$\Sigma$-sequence
$\mathfrak s$, put
$\lambda_{\mathfrak s}=\Omega^3(\delta_{\mathfrak s})$.
Then $\mathfrak s$ is in $\triangle_\theta$ precisely when
$\lambda_{\mathfrak s}=\lambda_N$ in the stable category. 
%Notice that
%$\lambda_N$ is not the inverse of $\theta_N$: the latter has source
%$\Sigma N$ and target $\Omega^{-3}(N)$.


%\subsection{Preprojective algebras}
%\label{sec:preprojective}


%The following realization is due to Dlab--Ringel.

%\begin{proposition}[{\cite[\S7, pp.~219--220]{DlabRingel1992}}]
%\label{prop:dlab-ringel}
%For $R_{n+1}=\kk[t]/(t^{n+1})$ and
%$G_n=\bigoplus_{i=1}^n R_{n+1}/(t^i)$, there %is %an algebra isomorphism
%\[
 %\Pi(A_n)\cong\ol{\End}_{R_{n+1}}(G_n).
%\]
%Here the quotient on the right is by morphisms factoring through the unique
%indecomposable projective $R_{n+1}$-module. With the right action
%$\phi\cdot a=\phi\circ a$, no opposite algebra %occurs.
%\end{proposition}

%\begin{proof}
%In \cite{DlabRingel1992}, the quotient of the Auslander algebra of $R_{n+1}$ by
%the ideal belonging to its projective summand is identified with the preprojective algebra
%of type $A_n$. That ideal is precisely the ideal of morphisms factoring
%through $R_{n+1}$, so their quotient is the stable endomorphism
%algebra of $G_n$.
%\end{proof}


%\begin{definition}\label{def:preprojective-algebra}
%Let $Q=(Q_0,Q_1,s,t)$ be a Dynkin quiver. Its %\emph{doubled quiver}
%$\overline Q$ has the same vertices and, for every arrow
%$a:s(a)\longrightarrow t(a)$ in $Q_1$, an additional reverse arrow
%$a^*:t(a)\longrightarrow s(a)$. Write $e_x$ for the stationary path at
%$x\in Q_0$. All paths are composed from right to left. %Over $\kk$, put
%\[
% \rho_Q=\sum_{a\in Q_1}
% \bigl(a^*\circ a-a\circ a^*\bigr)\in\kk\overline Q.
%\]
%The \emph{preprojective algebra} of $Q$ is
%\[
% \Pi(Q)=\kk\overline Q\Big/
% \bigl\langle e_x\rho_Qe_x\mid x\in Q_0\bigr\rangle,
%\]
%where the angle brackets denote the two-sided ideal generated by the listed
%vertex relations; see details in
%\cite[p.~158]{GelfandPonomarev1979}. We keep the conventional minus sign
%in $\rho_Q$, although it equals a plus sign over $\kk$.
%\end{definition}
\section{A pre-triangulated category which is not triangulated}
\label{sec:example}
The main result of this section is the proof of \Cref{thm:main} from the introduction, which is presented at the end of this section.
\subsection{The preprojective algebra \texorpdfstring{$\Pi(A_5)$}{Pi(A5)}}
\label{sec:algebra}

Let $\overline{A_5}$ be the following
doubled quiver of $A_5$: 
\[
\xymatrix@C=38pt{
1\ar@/^/[r]^{b_1}&
2\ar@/^/[l]^{a_1}\ar@/^/[r]^{b_2}&
3\ar@/^/[l]^{a_2}\ar@/^/[r]^{b_3}&
4\ar@/^/[l]^{a_3}\ar@/^/[r]^{b_4}&
5\ar@/^/[l]^{a_4}
}
\]
Let $\Lambda=\Pi(A_5)$ denote the preprojective algebra of type $A_5$ over $\mathbb F_2$, namely the quotient of the path algebra $\mathbb F_2 \overline{A_5}$ by the following relations:
\begin{equation}\label{eq:preproj-relations}
 a_1\circ b_1,\qquad
 b_i\circ a_i+a_{i+1}\circ b_{i+1}\ (1\le i\le3),\qquad
 b_4\circ a_4.
\end{equation}
See details in \cite[p.~158]{GelfandPonomarev1979}. Denote by $\modu \Lambda$ the category of finite-dimensional right $\Lambda$-modules. For $i\in\{1,2,3,4,5\}$, write $P_i=e_i\Lam$ for the corresponding indecomposable projective right $\Lam$-module.

Define the graph-reflection automorphism $\nu\in \operatorname{Aut}_{\kk}(\Lam)$
\begin{equation}\label{eq:nu}
 \nu(e_i)=e_{6-i},\qquad
 \nu(a_i)=b_{5-i},\qquad
 \nu(b_i)=a_{5-i}.
\end{equation}
It satisfies $\nu^2=\mathrm{Id}_\Lam$. For a right $\Lam$-module $N$, we write $N_\nu$ for the right module with twisted action
$n\mathbin{\cdot_\nu}a=n\cdot\nu(a)$. Indeed, if $N$ has dimension vector
\[
 \mathbf{dim}\,N=(t_1,t_2,t_3,t_4,t_5)
\]
and arrow matrices $N(a_i)$ and $N(b_i)$, $i=1,2,3,4$, then $N_\nu$ has dimension vector
\[
 \mathbf{dim}\,N=(t_5,t_4,t_3,t_2,t_1)
\]
and arrow matrices $N_\nu(a_i)=N(b_{5-i})$ and $N_\nu(b_i)=N(a_{5-i})$, $i=1,2,3,4$. We use the enveloping algebra
\[
 \Lam^e=\Lam^{\mathrm{op}}\otimes_{\kk}\Lam,
\]
and regard a $\Lam$-bimodule as a right $\Lam^e$-module via
$m\cdot(a^{\mathrm{op}}\otimes_{\kk} b)=a\circ m\circ b$.

%The following result is due to Brenner--Butler--King.

\begin{proposition}\label{prop:periodicity}
The algebra $\Lam$ is a $35$-dimensional self-injective algebra with Nakayama
automorphism $(-)_\nu$. There is an exact sequence of $\Lam$-bimodules
\begin{equation}\label{eq:bimodule-resolution}
0\longrightarrow\Lam_\nu\xlongrightarrow{j}Q_2
\xlongrightarrow{\partial_2}Q_1\xlongrightarrow{\partial_1}Q_0
\xlongrightarrow{\mu}\Lam\longrightarrow0,
\end{equation}
where 
\[
 Q_0=Q_2=\bigoplus_{i=1}^5\Lam e_i\otimes_{\kk} e_i\Lam
~~\text{ and }~~
 Q_1=\bigoplus_{x\in T}\Lam e_{t(x)}
 \otimes_{\kk} e_{s(x)}\Lam;
\]
here $T=\{a_i,b_i\mid 1\leq i\leq 4\}$. 
In particular,
\[
 \Omega^3_{\Lam^e}(\Lam)\cong\Lam_\nu
 \quad\text{and}\quad
 \Omega^3\cong(-)_\nu\cong\Omega^{-3}
 \quad\text{on }\ol{\modu}\Lam.
\]
The Auslander--Reiten translation $\tau$ equals to $\Omega^{-1}$. 
\end{proposition}

\begin{proof}
The Frobenius structure and the bimodule resolution of $\Lambda$ can be deduced by 
\cite[Theorems~4.8 and 4.9]{BrennerButlerKing2002}. It follows that $\Omega^3_{\Lam^e}(\Lam)\cong\Lam_\nu$. 

\vskip5pt

Let $N\in \modu \Lambda$. Note that the sequence \eqref{eq:bimodule-resolution} splits as an exact sequence of left $\Lambda$-modules. Thus, applying $N_\nu\otimes_\Lam-$ to \eqref{eq:bimodule-resolution} gives an exact sequence of right $\Lambda$-modules

\begin{equation}\label{eq:tensor-sequence}
0\to N\xrightarrow{J_N}P_N^0\xrightarrow{D_N^0}P_N^1
\xrightarrow{D_N^1}P_N^2\xrightarrow{Q_N}\Sigma N\to0, 
\end{equation}
where all $P_N^0, P_N^1$ and $P_N^2$ are projective-injective. This shows that $\Omega^3\cong(-)_\nu\cong\Omega^{-3}$ on $\ol{\modu}\Lam$. Finally, by \cite[Chapter~IV, Proposition~3.7] {AuslanderReitenSmalo1995}, $\tau\cong\Omega^2\circ(-)_\nu\cong\Omega^5\cong\Omega^{-1}$.
\end{proof}

Put $R=\kk[t]/(t^6)$, $V_i=R/(t^i)$, and $G=\bigoplus_{i=1}^5V_i$. A classical result due to Dlab and Ringel show the relationship between $\Lambda$ and $R$. 

\begin{proposition}\label{LambdaR}\cite[\S7, pp.~220--221]{DlabRingel1992}
There is a ring isomorphism
\[
 \Lam\cong\ol{\End}_R(G).
\]
The isomorphism is given by sending $a_i$ to the natural quotient $V_{i+1}\longrightarrow V_i$ and $b_i$ to the natural embedding $V_i\longrightarrow V_{i+1}$. 
\end{proposition}

\begin{remark}\label{conv:transported-triangulation}
Note that $\modu R$ is a Frobenius category, and $\ol{\modu}R$ is triangulated by \cite[Chapter~I, Theorem~2.6]{Happel1988}. In particular, the isomorphism
\[
 \varphi:\ol{\modu}R\xlongrightarrow{\sim}\mathcal P(\Lam).
\]
give a triangulated structure on $\mathcal F=P(\Lam)$. The corresponding suspension functor $\Sigma$ is just $(-)_\nu$. Denote the corresponding triangulation on $\cF$ by $\triangle_0$. By \Cref{thm:heller}, $\triangle_0$ induces a Heller comparison
\[
 \theta_0:\Sigma\xlongrightarrow{\sim}\Omega^{-3}.
\]
For $N\in\modu\Lam$, write
$(\lambda_0)_N=\Omega^3((\theta_0)_N)$.
\end{remark}


\subsection{A twist of comparison}\label{sec:parameter}

%We first isolate the general source of the twist.

For the definition and properties of Auslander--Reiten triangles, we refer readers to \cite[Chapter I. 4]{Happel1988}.

\begin{lemma}\label{lem:orbit-twist}
Let $\mathcal C$ be a Hom-finite Krull--Schmidt triangulated
$\kk$-category with Auslander--Reiten triangles and suspension $S$.
Let $\cO$ be a complete $S$-orbit of isomorphism classes of indecomposable
objects. Assume that for every representative $U\in\cO$, 
\[
 \End_{\mathcal C}(U)/\operatorname{rad}\End_{\mathcal C}(U)\cong\kk, 
\]
and suppose that there is a nonzero
$\Delta_U\in\operatorname{rad}\End_{\mathcal C}(U)$ such that
\begin{enumerate}
 \item $\Delta_U\circ f=0$ for every morphism
 $f:Y\longrightarrow U$ which is not a retraction;
 \item $g\circ\Delta_U=0$ for every morphism
 $g:U\longrightarrow Y$ which is not a section;
 \item $\Delta_{SU}=S(\Delta_U)$.
\end{enumerate}
For each $W\in\mathcal C$, we write $W=(\oplus_{i=1}^s W_1^i)\oplus (\oplus_{j=1}^tW_2^j)$ such that each $W_1^i$ and $W_2^j$ is indecomposable and $W_1^i\in\cO$ for each $1\leq i\leq s$, and  $W_2^j\notin \cO$ for each $1\leq j\leq t$, then we define $\Delta_W=(\oplus_{i=1}^s \Delta_{W_1^i})\oplus (\oplus_{j=1}^t 0_{W_2^j})$.
Then $\Delta_{\cO}$ is a natural transformation of the identity,
commutes with $S$, and satisfies $\Delta_{\cO}^2=0$. In particular, $\eps_{W}=\mathrm{Id}_W+\Delta_{W}$ is a unit which commutes with $S$ and
$\eps_{W}^{-1}=\eps_{W}$ for every $W\in\mathcal C$.
\end{lemma}

\begin{proof}
The hypotheses (1)(2) give naturality for morphisms between
nonisomorphic indecomposable objects. Thus, it remains to show that for each object $U\in \cO$ and each morphism $h\colon U\to U$, $\Delta_U\circ h=h\circ \Delta_U$.  Fix an object $U\in\cO$ and a morphism $h\in\End_{\mathcal C}(U)$. By the assumption that $\End_{\mathcal C}(U)/\operatorname{rad}\End_{\mathcal C}(U)\cong\kk$
 , there exists $c\in\mathbb F_2$ and $r\in\operatorname{rad}\End_{\mathcal C}(U)$ such that 
$
 h=c\,\mathrm{Id}_U+r.
$
Combining this with the assumptions (1)(2), we have
$\Delta_U\circ r=0=r\circ\Delta_U$, and hence
$$
\Delta_U\circ h=c\Delta_U=h\circ\Delta_U. $$
This proves the naturality on an
indecomposable object. Combining with the
Hom-finite Krull--Schmidt property of $\mathcal C$, we conclude that $\Delta_{\cO}$ is natural.


Because $\Delta_U\in \operatorname{rad}\End_{\mathcal C}(U)$ for each $U\in\cO$, it is neither a section nor a retraction, so
the hypotheses (1)(2) give $\Delta_U^2=0$. It follows that $\Delta_W^2=0$ for each $W\in\mathcal C$.  Combining with the definition of $\Delta_{\cO}$,  the action of $\Delta_{\cO}$ on the $S$-orbit in (3) yields that $\Delta_{\cO}$ commutes with
$S$. Finally,  note that
$\eps_{W}^2=\mathrm{Id}_W$ over $\kk$, we get $\eps_{W}^{-1}=\eps_{W}$.
\end{proof}

Next, the following module $M$ plays a key role in this article. 

\begin{definition}\label{def:module-M}
Let $M\in \modu\Lambda$ have dimension vector
\[
 \mathbf{dim}\,M=(2,2,2,1,0)
\]
and arrow matrices $x_i=M(a_i)$, $y_i=M(b_i)$
\[
\begin{aligned}
x_1&=\begin{pmatrix}0&1\\0&1\end{pmatrix},&
y_1&=\begin{pmatrix}1&1\\0&0\end{pmatrix},&
x_2&=\begin{pmatrix}1&0\\0&0\end{pmatrix},\\
y_2&=\begin{pmatrix}0&1\\0&0\end{pmatrix},&
x_3&=\begin{pmatrix}0&1\end{pmatrix},&
y_3&=\begin{pmatrix}1\\0\end{pmatrix},
\end{aligned}
\]
\end{definition}

Here are some properties of $M$. 

\begin{proposition}\label{prop:M}
The module $M$ is indecomposable, $\Omega(M)\cong M_\nu$ and $\Omega^2(M)\cong M$. Furthermore, 
\[
 \ol{\End}_\Lam(M) \cong\kk[\rho]/(\rho^2),
\]
\[
 \ol{\Hom}_\Lam(M,\Omega(M))=0
 =\ol{\Hom}_\Lam(\Omega(M),M).
\]
\end{proposition}

\begin{proof}
A direct computation shows that there is an exact sequence
\begin{equation}\label{PM}
     0\longrightarrow M_\nu
 \xlongrightarrow{\kappa}P_1\oplus P_3
 \xlongrightarrow{\pi}M
 \longrightarrow0.
\end{equation}

The components of $\kappa$ are
\[
 \kappa_1=0,\quad
 \kappa_2=
 \begin{pmatrix}
 0\\0\\1
 \end{pmatrix},\quad
 \kappa_3=
 \begin{pmatrix}
 0&1\\
 0&0\\
 0&1\\
 1&0
 \end{pmatrix},\quad
 \kappa_4=
 \begin{pmatrix}
 1&1\\
 0&0\\
 1&0
 \end{pmatrix},\quad
 \kappa_5=
 \begin{pmatrix}
 1&0\\
 0&1
 \end{pmatrix}.
\]
The components of the projective cover $\pi$ are
\[
 \pi_1=
 \begin{pmatrix}
 0&1\\
 1&0
 \end{pmatrix},\quad
 \pi_2=
 \begin{pmatrix}
 1&1&0\\
 1&0&0
 \end{pmatrix},\quad
 \pi_3=
 \begin{pmatrix}
 1&0&1&0\\
 0&1&0&0
 \end{pmatrix},\quad
 \pi_4=
 \begin{pmatrix}
 0&1&0
 \end{pmatrix},\quad
 \pi_5=0.
\]

Consequently, $\Omega(M)\cong M_\nu$. It follows that $\Omega^2(M)\cong(M_\nu)_\nu\cong M$. 

\vskip5pt

A direct computation shows that for each $h\in \End_\Lam(M)$, there are $c,s,t\in\kk$ such that
\[
 h_1=\begin{pmatrix}c&s\\0&c\end{pmatrix},\qquad
 h_2=c\,\mathrm{Id}_{\kk^2},\qquad
 h_3=\begin{pmatrix}c&t\\0&c\end{pmatrix},\qquad
 h_4=c\,\mathrm{Id}_{\kk},\qquad h_5=0, 
\]
and vice versa. Thus, $\End_\Lam(M)\cong \kk[u,v]/(u^2, uv, v^2)$. In particular, $\End_\Lam(M)$ is local and $M$ is indecomposable. And $\dim \operatorname{End}_\Lambda(M)=\dim \operatorname{End}_\Lambda(M_\nu)=\dim \operatorname{End}_\Lambda(\Omega(M))=3$. 

\vskip5pt

Similarly, a direct computation shows that for each $l\in \operatorname{Hom}_\Lambda(\Omega(M), M)$, $l_1,l_2,l_4,l_5=0$,
\[
 l_3=\begin{pmatrix}0&w\\0&0\end{pmatrix}, w\in \kk
\]
and vice versa. Thus, $\dim \operatorname{Hom}_\Lambda(\Omega(M), M)=\dim \operatorname{Hom}_\Lambda(\Omega(M), M)=1$. 

\vskip5pt

Now, applying $\operatorname{Hom}_\Lambda(-, M)$ to \eqref{PM}, one gets
\[
\begin{aligned}
\dim \ol{\operatorname{Hom}}_\Lambda(\Omega(M), M)&=\dim \operatorname{Hom}_\Lambda(M, M)-\dim \operatorname{Hom}_\Lambda(P_1\oplus P_3, M)+\dim \operatorname{Hom}_\Lambda(\Omega(M), M)\\
 &= 3-(2+2)+1\\
 &=0
\end{aligned}
\]
and $\dim \ol{\operatorname{Hom}}_\Lambda(M, \Omega(M))=\dim \ol{\operatorname{Hom}}_\Lambda(\Omega(M), M)=0$. Similarly, applying $\operatorname{Hom}_\Lambda(-, \Omega(M))$ to \eqref{PM}, one gets $\dim \ol{\operatorname{End}}_\Lambda(M)=\dim \ol{\operatorname{End}}_\Lambda(\Omega(M))=3-2+1=2$. In particular, 
\[
\ol{\End}_\Lam(M)\cong\kk[\rho]/(\rho^2). 
\]
This completes the proof. 
\end{proof}

Now, we want to apply \Cref{lem:orbit-twist} to $\ol{\modu} \Lambda$. It is well-known that $\ol{\modu} \Lambda$ has Auslander--Reiten triangles induced by Auslander--Reiten sequences in $\modu\Lambda$ and suspension $S= \Omega^{-1}$; see, for example, \cite{AuslanderReitenSmalo1995, Happel1988} or \cite[Theorem I.2.4]{ReitenVanDenBergh2002}. By \Cref{prop:M}, $\cO=\{M, \Omega(M)\}$ is a complete $S$-orbit, and $\ol{\End}_{\Lambda}(M)/\operatorname{rad}\ol{\End}_{\Lambda}(M)\cong\kk\cong \ol{\End}_{\Lambda}(\Omega(M))/\operatorname{rad}\ol{\End}_{\Lambda}(\Omega(M))$. Finally, let $\Delta_M=\rho$ be the unique nonzero non-isomorphism in $\ol{\End}_{\Lambda}(M)$ and $\Delta_{\Omega(M)}=\Omega(\rho)$. 

\begin{proposition}\label{prop:twisted-comparison}
There is a natural transformation
$\Delta_{\cO}:\mathrm{Id}_{\ol{\modu}\Lam}
\longrightarrow\mathrm{Id}_{\ol{\modu}\Lam}$ supported on $\cO=\{M,\Omega(M)\}$, 
such that $\Delta_{M}=\rho$, $\Delta_{\cO}^2=0$, and
$\Delta_{\cO}$ commutes with $S=\Omega^{-1}$. The components
$\eps_{U}=\mathrm{Id}_U+\Delta_{U}$ define a
suspension-commuting automorphism $\eps_{\cO}$, and
\begin{equation}\label{eq:twisted-comparison}
 (\theta_\eps)_N
 =(\eps_{\cO})_{\Omega^{-3}(N)}\circ(\theta_0)_N:
 \Sigma N\xlongrightarrow{\sim}\Omega^{-3}(N).
\end{equation}
defines an isomorphism of triangle functors
$\theta_\eps:\Sigma\xlongrightarrow{\sim}\Omega^{-3}$.
\end{proposition}

\begin{proof}
Note that there is an Auslander--Reiten triangle $ \Omega(M) \longrightarrow E\longrightarrow M \xlongrightarrow{\rho} M$ since $\tau M=\Omega^{-1}(M)=\Omega(M)$ and $\rho\in \ol{\End}_{\Lambda}(M)$ is the unique nonzero non-isomorphism. Thus, $\Delta_\cO$ satisfies all the conditions in \Cref{lem:orbit-twist}. This completes the proof. 
\end{proof}

\begin{corollary}\label{cor:pretriangulation}
The comparison $\theta_\eps$ in \eqref{eq:twisted-comparison} defines a pre-triangulation $\triangle_\eps$ on $(\cF,\Sigma)$.
\end{corollary}

\begin{proof}
By \Cref{prop:periodicity} and \Cref{conv:transported-triangulation}, $\Lam$ is self-injective and
$\Sigma=(-)_\nu$ is an auto-equivalence on $\cF=\mathcal P(\Lam)$. The desired result now follows from \Cref{thm:heller} and \Cref{prop:twisted-comparison}.
\end{proof}

\begin{lemma}\label{lem:twisted-lambda}
For every $N\in\modu\Lam$, applying $\Omega^3$ to
$(\theta_\eps)_N$ gives
\begin{equation}\label{eq:twisted-lambda}
 (\lambda_\eps)_N\coloneqq \Omega^3((\theta_\epsilon)_N)
 =(\mathrm{Id}_N+\Delta_{N})\circ(\lambda_0)_N.
\end{equation}
\end{lemma}

\begin{proof}
This follows from applying $\Omega^3$ to \eqref{eq:twisted-comparison}. Recall that $(\lambda_0)_N=\Omega^3((\theta_0)_N)$ by definition. 
\end{proof}

\subsection{Failure of the four-by-four property}
\label{sec:square}
Next, our aim is to describe the obstruction of $\mathbf{TR4}$ for $\triangle_\eps$.  

In the following, we construct the modules $A,B$ based on $M$ via the Auslander-- Reiten theory. See details in \cite{AuslanderReitenSmalo1995}.
\begin{definition}\label{def:modules-AB}
By the stable isomorphism of $M$ in \Cref{prop:M}, there exists a unique Auslander--Reiten sequence in $\operatorname{Ext}^1_\Lambda(\Omega(M),M)$ up to isomorphism. 
Let
\begin{equation}\label{eq:AR-sequence}
 0\longrightarrow M\xlongrightarrow{\iota}A
 \xlongrightarrow{\pi}\Omega(M)\longrightarrow0
\end{equation}
be a representative of the Auslander--Reiten sequence in $\operatorname{Ext}^1_\Lambda(\Omega(M),M)$, and define $A$ to be its
middle term. Choose a short exact sequence
\begin{equation}\label{eq:extension}
 0\longrightarrow A\xlongrightarrow{u}B
 \xlongrightarrow{p}M\longrightarrow0
\end{equation}
representing the nonzero element of $\operatorname{Ext}^1_\Lam(M,A)$, and
define $B$ to be its middle term; the existence and uniqueness of the nonzero class follows from the isomorphism $\operatorname{Ext}^1_\Lambda(M,A)\cong \ol{\Hom}_\Lam(\Omega(M),A)$ and \Cref{prop:intrinsic-relations}.
\end{definition}

\begin{proposition}\label{prop:intrinsic-relations}
%The second sequence in \Cref{def:modules-AB} exists and, up to equivalence,
%is the unique nonsplit extension of $M$ by $A$. 
In $\ol{\modu}\Lam$ one has $\ol{\Hom}_\Lam(\Omega(M),A)\cong\kk$, 
$
 S^2A\cong A,~S^2B\cong B.
$
Also, there are isomorphisms
\[\ol{\Hom}_\Lam(\Omega(M),B)=0=\ol{\Hom}_\Lam(B,\Omega(M)),
\]
and
\begin{equation}\label{eq:endomorphism-algebras}
 \ol{\End}_\Lam(A)\cong\kk[x]/(x^2),\qquad
 \ol{\End}_\Lam(B)\cong\kk[x,y]/(x^2,y^2).
\end{equation}
In particular, $A$ and $B$ are indecomposable in the stable category, and
neither has a direct summand in the orbit $\{M,\Omega(M)\}$. The period-two
isomorphisms may be chosen compatibly with the distinguished triangles induced by
\eqref{eq:AR-sequence} and \eqref{eq:extension}.
\end{proposition}

\begin{proof}
By \Cref{prop:M}, $M\cong \Omega^2(M)$. Under this identification, the distinguished triangle in $\ol{\modu}\Lam$ corresponding 
\eqref{eq:AR-sequence} has the form
\begin{equation}\label{eq:AR-triangle-Y}
 M\xlongrightarrow{\iota}A\xlongrightarrow{\pi}\Omega(M)
 \xlongrightarrow{\rho_\Omega(M)}\Omega(M),
\end{equation}
where $\rho_\Omega(M)=S(\rho)$ is the nonzero radical element of
$\ol{\End}_\Lam(\Omega(M))\cong\kk[\rho_\Omega(M)]/(\rho_\Omega(M)^2)$. Applying the 
functor $\ol{\Hom}_\Lam(\Omega(M),-)$ to \eqref{eq:AR-triangle-Y}, we could get a long exact sequence of abelian groups. Combining with 
$
 \ol{\Hom}_\Lam(\Omega(M),M)=0
$
from \Cref{prop:M}, we get that $ \ol{\Hom}_\Lam(\Omega(M),A)=\kk f$ and $\pi f=\rho_\Omega(M)$ for some $f\in\ol{\Hom}_\Lambda(\Omega(M),A)$. Similarly, applying the 
functor $\ol{\Hom}_\Lambda(M,-)$ and $\ol{\Hom}_\Lam(-,\Omega(M)$ to \eqref{eq:AR-triangle-Y},  we could conclude by \Cref{prop:M} that
\[
\begin{aligned}
 \ol{\Hom}_\Lam(M,A)&=\kk\iota,&
 \ol{\Hom}_\Lam(A,M)&=\kk a,\\
 \ol{\Hom}_\Lam(A,\Omega(M))&=\kk\pi,& a\iota=\rho
&\end{aligned}
\]
for some  $a$. Combining with these equations, by applying $\ol{\Hom}_\Lambda(-,A)$ to \eqref{eq:AR-triangle-Y}, we get 
\[
 \ol{\End}_\Lam(A)
 =\kk\,\mathrm{Id}_A\oplus\kk n,
 \qquad n=\iota a=f\pi,\qquad n^2=0.
\]
This proves the first isomorphism in \eqref{eq:endomorphism-algebras}.
Applying $S^2$ to \eqref{eq:AR-triangle-Y} and using $S^2M\cong M$ and
$S^2\Omega(M)\cong \Omega(M)$, uniqueness of Auslander--Reiten triangles gives
$S^2A\cong A$.

Next,
consider the isomorphism $
 \operatorname{Ext}^1_\Lam(M,A)
\cong\ol{\Hom}_\Lam(\Omega(M),A)$. This is isomorphic to $\kk$.
Thus the nonzero extension \eqref{eq:extension} exists. Let $\eta:M\longrightarrow SA$ be its connecting morphism.
We may take $\eta=Sf$. In the stable category $\ol{\modu}\Lam$
$(Su)\eta=0$ (see \Cref{rem:consecutive-zero}), and hence $uf=0$. Applying
$\ol{\Hom}_\Lam(\Omega(M),-)$ to this triangle and using
$\ol{\Hom}_\Lam(\Omega(M),M)=0$ yields
$\ol{\Hom}_\Lam(\Omega(M),B)=0$. Dually, applying
$\ol{\Hom}_\Lam(-,\Omega(M))$ and using $\pi f=\rho_\Omega(M)\ne0$ gives
$\ol{\Hom}_\Lam(B,\Omega(M))=0$.

It remains only to record the endomorphisms of the second middle term.
By applying stable Hom to the two triangles above and the relations $\rho^2=\rho_\Omega(M)^2=0$, we conclude that
$ \dim_{\kk}\ol{\End}_\Lam(B)=4.
$
More precisely, its radical has a basis $x,y,xy$, with $x^2=y^2=0$ and $xy=yx$,
which proves the second isomorphism in
\eqref{eq:endomorphism-algebras}. 

Both endomorphism algebras in \eqref{eq:endomorphism-algebras} are local.
Moreover, $\ol{\Hom}_\Lam(\Omega(M),A)$ has dimension one, whereas
$\ol{\Hom}_\Lam(\Omega(M),M)$ and $\ol{\End}_\Lam(\Omega(M))$ have dimensions zero and two,
respectively. Thus $A$ is not isomorphic to $M$ or $\Omega(M)$. The stable
endomorphism algebra of $B$ has dimension four, while those of $M$ and $\Omega(M)$
have dimension two, so $B$ is not isomorphic to either of them. This proves
the assertion about the orbit.

Finally, $S^2$ preserves the unique nonzero class in
$\operatorname{Ext}^1_\Lam(M,A)$ and fixes its end terms up to the chosen
period-two isomorphisms. Hence $S^2B\cong B$. The uniqueness in both
constructions allows the period-two isomorphisms to be chosen compatibly
with the two distinguished triangles.
\end{proof}

\begin{lemma}\label{lem:complete-horseshoe}
There are distinguished exact $3$-$\Sigma$-sequences for $\triangle_0$
\[
 T_X:\quad
 E_{X,0}\xlongrightarrow{d_{X,0}}E_{X,1}
 \xlongrightarrow{d_{X,1}}E_{X,2}
 \xlongrightarrow{d_{X,2}}\Sigma E_{X,0},
 \qquad X\in\{A,B,M\},
\]
and a degreewise split short exact sequence of $\Sigma$-periodic complexes
\begin{equation}\label{eq:complete-horseshoe}
 0\longrightarrow T_A\xlongrightarrow{x_\bullet}T_B
 \xlongrightarrow{q_\bullet}T_M\longrightarrow0
\end{equation}
whose sequence of first kernels is \eqref{eq:extension}. In particular,
$x_i:E_{A,i}\longrightarrow E_{B,i}$ is a split monomorphism with cokernel
$E_{M,i}$ for $0\leq i\leq2$.
\end{lemma}

\begin{proof}
Choose complete projective-injective resolutions of $A$ and $M$ and apply
the horseshoe lemma to \eqref{eq:extension}. We get a degreewise split
short exact sequence of complete resolutions whose sequence of zeroth
cycles is \eqref{eq:extension}. Its third cosyzygies represent the triangle
obtained by applying $S^3$ to the stable triangle of
\eqref{eq:extension}. Since
$\theta_0:\Sigma\longrightarrow S^3$ is an isomorphism of triangle
functors, its components on $A,B,M$ identify this triangle with the
$\Sigma$-suspension of the original one. The comparison theorem lifts this
identification to the three complete resolutions. By adjoining
projective-injective summands to form exact $3$-$\Sigma$-sequences and taking truncations, we get
\eqref{eq:complete-horseshoe}, and every row has comparison
$\theta_0$.
\end{proof}

The rows $T_A$ and $T_B$ also belong to $\triangle_\eps$. Indeed,
\Cref{prop:intrinsic-relations} shows that neither $A$ nor $B$ has a stable
direct summand in $\{M,\Omega(M)\}$, so their comparisons are unchanged by the
twist. The first two component sequences in \eqref{eq:complete-horseshoe}
give split distinguished columns
\[
 E_{A,i}\xlongrightarrow{x_i}E_{B,i}
 \xlongrightarrow{q_i}E_{M,i}\longrightarrow\Sigma E_{A,i}
\]
for each $i=0,1$.

\begin{lemma}\label{lem:relative-filler}
Keep the first two components $x_0,x_1$ of $x_\bullet$ fixed. If
$x:E_{A,2}\longrightarrow E_{B,2}$ is any third component which completes
them to a morphism $T_A\longrightarrow T_B$, then
$
 x=sx_2
$
for an automorphism
\[
 \bigl(\mathrm{Id}_{E_{B,0}},\mathrm{Id}_{E_{B,1}},s,
 \Sigma\mathrm{Id}_{E_{B,0}}\bigr)
\]
of $T_B$. In particular, every such $x$ is a split monomorphism.
\end{lemma}

\begin{proof}
Set$
 C_X=\operatorname{Coker}d_{X,1}$ for each $X\in\{A,B\}$ and  $Z=\operatorname{Ker}d_{B,2}=\operatorname{Im}d_{B,1}.
$
Take $v=x-x_2$. Then we have
$
 vd_{A,1}=0,~d_{B,2}v=0,
$
and hence corresponds to a map $\bar v:C_A\longrightarrow Z$. That is, $v$ is the factorization  $
E_{A,2}\twoheadrightarrow C_A\xlongrightarrow{\bar n}
 Z\hookrightarrow E_{B,2}.
 $
Similarly, an
endomorphism $n$ of $E_{B,2}$ satisfying
$
 nd_{B,1}=0,~ d_{B,2}n=0
$
corresponds to a map $\bar n:C_B\longrightarrow Z$. Under these
identifications the map $n\mapsto nx_2$ is the restriction map
\begin{equation}\label{eq:filler-restriction}
 \Hom_\Lam(C_B,Z)\longrightarrow\Hom_\Lam(C_A,Z),
 \qquad \bar n\longmapsto\bar n\bar x_2,
\end{equation}
where $\bar x_2:C_A\longrightarrow C_B$ is induced by $x_2$.

Next, we prove that \eqref{eq:filler-restriction} is surjective. Note that in the stable category,
$C_X\simeq\Sigma X\simeq S^3X$ and $Z\simeq S^2B.
$
Under the identifications in
\Cref{prop:intrinsic-relations}, the stable form of
\eqref{eq:filler-restriction} is therefore
\begin{equation}\label{eq:stable-filler-restriction}
 (Su)^*:\ol{\Hom}_\Lam(SB,B)\longrightarrow
 \ol{\Hom}_\Lam(SA,B).
\end{equation}

Let $\eta:M\longrightarrow SA$ be the connecting morphism of
\eqref{eq:extension}. Suspending the Auslander--Reiten triangle
\eqref{eq:AR-triangle-Y} gives
\begin{equation}\label{eq:AR-triangle-M}
 \Omega(M)\longrightarrow SA\xlongrightarrow{q}M
 \xlongrightarrow{\rho}M.
\end{equation}
The map $q$ is right almost split. Since $\rho$ factors through $q$ and
$\ol{\Hom}_\Lam(M,SA)=\kk\eta$, one has
\begin{equation}\label{eq:q-eta}
 q\eta=\rho.
\end{equation}
For any $\bar h:SA\longrightarrow B$, the equality
$\ol{\Hom}_\Lam(\Omega(M),B)=0$ gives $\bar h=\ell q$ for some
$\ell:M\longrightarrow B$. The map $\ell$ is not a section, because $B$
has no stable direct summand isomorphic to $M$. Since $\rho$ is almost
vanishing, $\ell\rho=0$, and hence
$
 \bar h\eta=\ell q\eta=\ell\rho=0.
$
Applying $\ol{\Hom}_\Lam(-,B)$ to the suspended triangle of
\eqref{eq:extension} now shows that
\eqref{eq:stable-filler-restriction} is surjective.

Consequently, for any $\bar v:C_A\longrightarrow Z$ there is
$\bar n:C_B\longrightarrow Z$ such that
$\bar v-\bar n\bar x_2$ factors through a projective module $P$. The
algebra $\Lam$ is self-injective, so $P$ is injective and the factor
$C_A\longrightarrow P$ extends across the monomorphism
$\bar x_2:C_A\longrightarrow C_B$. This yields that the restriction map
\eqref{eq:filler-restriction} is surjective.

Choose $n$ with $v=nx_2$. Its factorization has the form
\[
 E_{B,2}\twoheadrightarrow C_B\xlongrightarrow{\bar n}
 Z\hookrightarrow E_{B,2}.
\]
Since $Z=\operatorname{Im}d_{B,1}$, the quotient
$E_{B,2}\twoheadrightarrow C_B$
vanishes on $Z$. It follows that $n^2=0$. Thus
$s=\mathrm{Id}_{E_{B,2}}+n$ is an automorphism. Moreover,
$sd_{B,1}=d_{B,1}$ and $d_{B,2}s=d_{B,2}$. This yields that the desired
strict automorphism of $T_B$, and $x=sx_2$.
\end{proof}

\begin{theorem}\label{thm:fixed-boundary-failure}
The pre-triangulated category
$(\mathcal{P}(\Lambda),\Sigma,\triangle_\eps)$ does not have the fixed-boundary
four-by-four property.
\end{theorem}

\begin{proof}
Take the rows $T_A,T_B$ and the first two split columns
as above. Note that these triangles are distinguished in $\triangle_\eps$.

Suppose that the above choice admits a four-by-four completion. Its third
vertical arrow $x:E_{A,2}\longrightarrow E_{B,2}$ completes
$x_0,x_1$ to a morphism of the two rows. By
\Cref{lem:relative-filler}, an automorphism of $T_B$  transports $x$ to $x_2$. This does not
change the prescribed column, so we may assume that the third vertical
arrow is $x_2$.

The map $x_2$ is a split monomorphism. A distinguished triangle beginning
with a split monomorphism is isomorphic to its split cokernel triangle. We may therefore assume
the third column is
\[
 E_{A,2}\xlongrightarrow{x_2}E_{B,2}
 \xlongrightarrow{q_2}E_{M,2}\longrightarrow\Sigma E_{A,2}.
\]
Since every $q_i$ is an epimorphism, commutativity uniquely determines all
three arrows in the third row. They are precisely the quotient arrows in
$T_M$; see \Cref{lem:complete-horseshoe}. Hence the third row of the four-by-four completion is $T_M$.

Since $T_M$ is in $\triangle_0$, the complete comparison of $T_M$ is $\theta_{0,M}$, or,
after applying $\Omega^3$, $(\lambda_0)_M$.  By above $T_M$ is in $\triangle_\eps$. It follows that
$
 (\lambda_\eps)_M=(\mathrm{Id}_M+\rho)(\lambda_0)_M;
$
see \Cref{lem:twisted-lambda}. Combining with that $T_M$ is in both $\triangle_0$ and $\triangle_\eps$, we conclude by \Cref{thm:heller} that $(\lambda_\eps)_M=(\lambda_0)_M$, and hence $(\mathrm{Id}_M+\rho)(\lambda_0)_M=(\lambda_0)_M$. This is a contradiction 
because $(\lambda_0)_M$ is an
isomorphism and $\rho\ne0$. This completes the proof. 
\end{proof}

We are now ready to prove the main result of this article.

{\bf Proof of \Cref{thm:main}.} By \Cref{cor:pretriangulation},
$\triangle_\eps$ satisfies $\mathbf{TR1}$--$\mathbf{TR3}$. If it satisfied
$\mathbf{TR4}$, then \Cref{lem:tr4-implies-fixed-boundary} would give the
four-by-four property, contrary to
\Cref{thm:fixed-boundary-failure}. Hence $\mathbf{TR4}$ fails.

\begin{ack}
    The mathematical search leading to the construction was carried out with substantial assistance from Eureka, an autonomous multi-agent mathematical-reasoning system. We also used GPT-5.6 Sol to check and revise this manuscript. We would like to thank the JIUCHONG team at the University of Science and Technology of China for providing us with access to Eureka. We are grateful to Jie Li and Tianyang Sun for their assistance in using Eureka. The third author thanks Bernhard Keller and Amnon Neeman, separately, for helpful private discussions and for patiently answering his question about whether a pre-triangulated category is triangulated.

Xiao-Wu Chen was supported by 
National Key R\textup{\&}D Program of China (No. 2024YFA1013801), and
the National Natural Science Foundation of China (No.s 12325101 and 12131015). Jian Liu was supported by the National Natural Science Foundation of China (No. 12401046). Xue-Song Lu and Chencheng Zhang were supported by National Natural Science Foundation of China (No. 12131015).
\end{ack}

\begin{thebibliography}{GKO13}

\bibitem[Ami07]{Ami07} C. Amiot, On the structure of triangulated categories with finitely many indecomposables, Bull. Soc. Math. France 135 (2007), no. 3, 435–474.

\bibitem[ARS95]{AuslanderReitenSmalo1995}
M.~Auslander, I.~Reiten, and S.~O.~Smal{\o},
\emph{Representation Theory of Artin Algebras},
Cambridge Studies in Advanced Mathematics, vol.~36,
Cambridge University Press, Cambridge, 1995.

\bibitem[Bel00]{Beligiannis2000}
A.~Beligiannis,
\emph{Relative homological algebra and purity in triangulated categories},
J. Algebra \textbf{227} (2000), no.~1, 268--361.

\bibitem[BBK02]{BrennerButlerKing2002}
S.~Brenner, M.~C.~R.~Butler, and A.~D.~King,
\emph{Periodic algebras which are almost Koszul},
Algebr. Represent. Theory \textbf{5} (2002), no.~4, 331--368.

\bibitem[DR92]{DlabRingel1992}
V.~Dlab and C.~M.~Ringel,
\emph{The module theoretical approach to quasi-hereditary algebras},
in \emph{Representations of Algebras and Related Topics},
London Math. Soc. Lecture Note Ser., vol.~168,
Cambridge University Press, Cambridge, 1992, pp.~200--224.

\bibitem[GKO13]{GeissKellerOppermann2013}
C.~Geiss, B.~Keller, and S.~Oppermann,
\emph{$n$-angulated categories},
J. Reine Angew. Math. \textbf{675} (2013), 101--120.

\bibitem[GP79]{GelfandPonomarev1979}
I.~M.~Gel'fand and V.~A.~Ponomarev,
\emph{Model algebras and representations of graphs},
Funct. Anal. Appl. \textbf{13} (1979), no.~3, 157--166.




\bibitem[Hap88]{Happel1988}
D.~Happel,
\emph{Triangulated Categories in the Representation Theory of
Finite-Dimensional Algebras},
London Math. Soc. Lecture Note Ser., vol.~119,
Cambridge University Press, Cambridge, 1988.


\bibitem[Har66]{Hartshorne1966}
R. Hartshorne,
\emph{Residues and Duality},
 Lecture Notes in Math., vol.~20,
Springer, Berlin, 1966.

\bibitem[Hel68]{Heller1968}
A.~Heller,
\emph{Stable homotopy categories},
Bull. Amer. Math. Soc. \textbf{74} (1968), no.~1, 28--63.



\bibitem[IVdB02]{ReitenVanDenBergh2002}
I.~Reiten and M.~Van den Bergh,
\emph{Noetherian hereditary abelian categories satisfying Serre duality},
J. Amer. Math. Soc. \textbf{15} (2002), no.~2, 295--366.

\bibitem[May01]{May2001}
J.~P.~May,
\emph{The additivity of traces in triangulated categories},
Adv. Math. \textbf{163} (2001), no.~1, 34--73.

\bibitem[Nee01]{Neeman2001}
A.~Neeman,
\emph{Triangulated Categories},
Annals of Mathematics Studies, vol.~148,
Princeton University Press, Princeton, NJ, 2001.

\bibitem[Pup67]{Puppe1967}
D.~Puppe,
\emph{Stabile Homotopietheorie~I},
Math. Ann. \textbf{169} (1967), no.~2, 243--274.

\bibitem[Ver77]{Verdier1977}
J.-L.~Verdier,
\emph{Cat\'egories d\'eriv\'ees},
in \emph{SGA $4\frac12$}, Lecture Notes in Math., vol.~569,
Springer, Berlin, 1977, pp.~262--311.

\end{thebibliography}

\end{document}
